\documentclass[11pt,a4paper]{article}

\usepackage[T1]{fontenc}
\usepackage[utf8]{inputenc}
\usepackage{lmodern}
\usepackage[a4paper,margin=2.6cm]{geometry}
\usepackage{microtype}
\usepackage{amsmath,amssymb,amsthm,mathtools}
\usepackage[hidelinks]{hyperref}

\newtheorem{theorem}{Theorem}[section]
\newtheorem{corollary}[theorem]{Corollary}
\theoremstyle{definition}
\newtheorem{definition}[theorem]{Definition}
\newtheorem{example}[theorem]{Example}

\numberwithin{equation}{section}

\title{Organizational Depth at Finite Time:\\
A Fixed-Resolution No-Go Boundary}
\author{}
\date{}

\begin{document}
\maketitle

\begin{abstract}
When does infinite organizational depth represent a physically meaningful
finite-time escape rather than increasingly fine indexing of a finite
transformation? We reformulate this question within Organizational
Accessibility around three results. First, a sequential stochastic
accessibility process with transition rate
$\lambda_n=N_n\nu_np_n$ is non-explosive when total material is bounded,
each independently counted realization requires a positive minimum
material allocation, and candidate generation per realization is uniformly
bounded. Second, a continuum cascade can retain bounded aggregate
throughput while its specific turnover diverges because the inventory
assigned to successive levels tends to zero. This identifies a
shrinking-scale route but does not derive its accelerating kinetics
physically.

The central mathematical result is a conditional fixed-resolution no-go
theorem. If transition costs and durations satisfy
$\varepsilon_n\tau_n\ge c_n d_{\rm phys,n}^2$, finite total charged energy
and time imply
$\sum_n\sqrt{c_n}\,d_{\rm phys,n}<\infty$.
A uniform positive kinetic coefficient therefore bounds accumulated
physical distance, and a uniform physical-to-operational comparison bounds
the number of transitions at any fixed operational resolution.

For finite-state continuous-time Markov jump dynamics, the metric comparison
can be discharged rather than assumed. Total variation is the distance in an
activity--entropy speed limit and contracts under every stochastic measurement
channel. If $\Sigma_*$ is total entropy production and $\mathcal N_*$ the
expected total number of jumps, then
\[
\sum_n d_{\rm op,n}\le\sqrt{\Sigma_*\mathcal N_*/2}.
\]
The associated consecutive-transition count is asymptotically sharp: a
near-equilibrium driven edge approaches equality. But a two-state shuttle can
saturate that count while revisiting only two states, so consecutive
fixed-resolution transitions are not themselves organizational depth.

We therefore define the retained packing depth $D_\delta$ as the largest
number of visited retained states that are mutually separated by at least
$\delta$ under one fixed operational metric. Chronologically ordering such a
set creates $D_\delta-1$ disjoint gaps, each of operational distance at least
$\delta$, and hence
\[
D_\delta
\le
1+\frac1\delta\sqrt{\frac{\Sigma_*\mathcal N_*}{2}}.
\]
Independently,
$D_\delta$ is bounded by the $\delta$-packing number of the accessible
operational state space. For a finite $m$-state probability simplex,
largest-remainder lattice rounding gives
\[
D_\delta
\le
\binom{k+m-1}{m-1},
\qquad
k=\left\lfloor\frac{m}{2\delta}\right\rfloor+1.
\]
Thus finite-state retained depth is constrained both by dynamical resources
and by geometric repertoire, and both factors can become the active bound. A bounded-power refinement example still shows
why infinite labels can accumulate when resolved scale shrinks.

These are conditional physical boundaries, not a universal prohibition of
singular dynamics. The remaining challenge is constructive: any proposed
finite-time infinite-depth system must expose how it creates unbounded
operational packing capacity, evades the dynamical-resource bound, or leaves
the stated model class.
\end{abstract}


\section{Introduction: from singularity to physical distinguishability}

Organizational Accessibility asks how currently realized organization
changes which successor organizations can be generated, established,
and retained \cite{OA8}. Its central historical claim is conditional:
previously successful organization can become infrastructure for
subsequent organization. This does not imply that organizational
change accelerates indefinitely, or that increasing organizational
depth corresponds to increasing physical change.

A proposed finite-time hierarchy introduces a further question.
Suppose retained states are labelled
\[
x_0,x_1,x_2,\ldots
\]
and their transition times accumulate before a finite deadline.
What makes these labels represent distinct physical organizations?
Without an answer, an infinite hierarchy may describe a finite
trajectory divided into increasingly small pieces.

The present paper therefore asks:
\begin{quote}
Under what conditions can infinite organizational depth remain
physically meaningful at finite time, finite charged energy, and fixed
operational resolution?
\end{quote}

The analysis separates three issues. A stochastic event model requires
a sufficient rate of completed transitions. A continuum throughput model
requires an interpretation of increasingly rapid turnover of shrinking
quantities. An organizational interpretation additionally requires
physically realizable distinctions between successive retained states.

The general boundary is conditional. Finite time and energy exclude
infinitely many fixed-resolution transitions when an appropriate
speed--distance inequality holds uniformly and its physical metric
controls operational distinguishability. Neither condition follows
merely from assigning an increasing organizational level number.

For one important model class, however, the metric comparison is not an
extra assumption. In finite-state Markov jump dynamics the
Shiraishi--Funo--Saito speed limit is expressed directly in total
variation, while total variation contracts under any stochastic
measurement channel. This gives a model-specific operational path-length
theorem. A subsequent tightness check shows that the resulting
consecutive-transition count is asymptotically sharp but is not by itself a
depth measure: repeated shuttling between two states can attain it. The paper
therefore distinguishes transition count from retained packing depth and
combines the thermodynamic path-length bound with a geometric packing bound.
In this setting the abstract kinetic coefficient can be identified with
inverse average activity.

The generic mathematical ingredients are established: pure-birth
non-explosion, Cauchy--Schwarz inequalities, stochastic thermodynamic
speed limits, optimal-transport bounds, and metric convergence
\cite{Feller1971,Shiraishi2018,Aurell2011,VuSaito2023,Nagayama2025,Zhang2001}.
The proposed contribution is their application to organizational depth,
including explicit closure of one stochastic route, separation of
aggregate throughput from specific turnover, a fixed-resolution
criterion, and identification of dynamical activity as the concrete
escape resource in a finite-state stochastic class.

\section{Organization as amount, configuration, and operational distinction}

Three descriptions of organization must be distinguished.

\emph{Organization as amount} is an extensive quantity whose increase
requires an increase in a material or energetic stock. The component
abundances in the resource-limited model of Organizational Accessibility
have this interpretation.

\emph{Organization as configuration or history} concerns arrangement,
architecture, or retained causal structure. Such quantities need not
increase in direct proportion to material stock. The same physical
substrate may participate in many successive configurations.

\emph{Organization as an operational distinction} concerns differences
between retained states that can be resolved under a specified
measurement standard. Different labels do not establish such a
distinction.

The physical state description must include the subsystems needed to
implement the claimed transformation. These can include processors,
catalysts, controllers, memory, interfaces, fuel stores, and retained
records. Omitting them can turn a physical cost into an apparently
cost-free external condition.

\subsection*{Transitions and accounting}

Set $t_0=0$ and write
\[
t_{n+1}=t_n+\tau_n,
\qquad
\tau_n>0.
\]
A finite-time sequential accumulation has
\begin{equation}
T_*=\sum_{n=0}^{\infty}\tau_n<\infty.
\label{eq:time-sum}
\end{equation}

Let $\varepsilon_n\ge0$ be a \emph{charged energetic cost}. This means
a nonnegative expenditure that can legitimately be counted against
a cumulative budget:
\begin{equation}
\sum_{n:\,t_{n+1}\le T}\varepsilon_n
\le B(T)<\infty
\qquad\text{for every finite }T.
\label{eq:budget}
\end{equation}

Raw signed work, heat, stored energy, and excess dissipation are not
interchangeable choices for $\varepsilon_n$. The selected cost must
match both the budget and the speed--distance inequality used below.

For example, a model may justify
\[
B(T)=E_0+\int_0^T P_{\rm av}(t)\,dt
\]
after accounting for finite initial stores and the usable free-energy
sources that pay the charges. Bounded incoming power alone does not
bound gross internal processing power, nor does it justify charging
the same recoverable energy repeatedly as if it were irreversibly
consumed.

The finite-action results below are statements about sequences
satisfying explicit inequalities. They apply pathwise when their
premises hold pathwise. They can also apply to prescribed sequences
of thermodynamic protocols with deterministic durations and
ensemble-mean costs, provided the same ensemble interpretation is
used throughout. They do not convert an average thermodynamic
inequality into an almost-sure event-cost bound.

\subsection*{Material concentration is a separate question}

The normalized finite-dimensional production model in
Organizational Accessibility has a material balance of the form
$\dot M=J-M$. Its abundance bounds concern organization as amount.
Spatial extensions raise a different regularity question:
control of integrated material does not automatically control
pointwise concentration or stronger norms
\cite{PierreSchmitt1997}.

The present fixed-resolution results neither establish nor exclude
all such spatial singularities. They concern specified sequences of
retained transformations, their physical costs, and their operational
interpretation.


\section{Closure of stochastic route C}

Consider a sequential pure-birth accessibility process
\[
O_0\longrightarrow O_1\longrightarrow O_2\longrightarrow\cdots.
\]
At level $n$, let
\begin{equation}
\lambda_n=N_n\nu_np_n,
\label{eq:lambda}
\end{equation}
where $N_n$ is the number of active searching realizations,
$\nu_n$ is their candidate-generation rate per realization, and
$p_n$ is the probability that a candidate belongs to the specified
successor class and establishes.

Because $p_n$ is a probability,
\[
0\le p_n\le1.
\]
This bound is part of the model definition, not an additional
empirical hypothesis.

\begin{theorem}[Closure of stochastic accessibility]
\label{thm:C}
Suppose, at every reachable level,
\[
M_n\le M_*<\infty,
\qquad
M_n\ge m_{\min}N_n,
\qquad
m_{\min}>0,
\]
and
\[
0\le\nu_n\le\nu_{\max}<\infty.
\]
Then the pure-birth process with rates \eqref{eq:lambda} is
non-explosive.
\end{theorem}

\begin{proof}
The material inequalities give
\[
N_n\le\frac{M_n}{m_{\min}}
\le\frac{M_*}{m_{\min}}.
\]
Therefore
\begin{equation}
\lambda_n\le
\Lambda:=
\frac{M_*}{m_{\min}}\nu_{\max}<\infty.
\label{eq:lambda-bound}
\end{equation}

If $\Lambda=0$, no transition occurs. If a reachable
$\lambda_n=0$, the sequential process stops there.

Otherwise,
\[
\sum_{n=0}^{\infty}\frac1{\lambda_n}
\ge
\sum_{n=0}^{\infty}\frac1\Lambda
=\infty.
\]
This meets the classical non-explosion criterion for the pure-birth
model \cite{Feller1971}.

A direct almost-sure argument is also available. Represent the
independent holding times as
\[
X_n=\frac{E_n}{\lambda_n},
\qquad
E_n\overset{\rm iid}{\sim}{\rm Exp}(1).
\]
By \eqref{eq:lambda-bound},
\[
X_n\ge\frac{E_n}{\Lambda}.
\]
The strong law gives $\sum_nE_n=\infty$ almost surely, and hence
\[
\sum_nX_n=\infty
\qquad\text{almost surely}.
\]
\end{proof}

The additive processor inequality
$M_n\ge m_{\min}N_n$ is substantive. It is not valid if several
overlapping descriptions of the same physical machinery are counted
as independent full-rate realizations.

For the reciprocal-rate series to converge, it is necessary that
$\lambda_n\to\infty$. Since $p_n\le1$, this requires
$N_n\nu_n\to\infty$. Thus a C-type explosion must abandon at least
one uniform bound or accounting assumption used in
Theorem~\ref{thm:C}: bounded material, a positive minimum allocation
per realization, bounded generation rate per realization, or additive
processor accounting. In parameterized limits these failures may
appear as
\[
M_*\to\infty,\qquad
m_{\min}\downarrow0,\qquad
\nu_{\max}\to\infty.
\]
Alternatively, a model may depart from the sequential one-step event
description, for example by allowing an unbounded depth increment
within one event.

Failure of an assumption is not sufficient for explosion. Unbounded
rates can still have a divergent reciprocal series.

Endogenizing $N_n,\nu_n,p_n$ remains relevant to finite accessibility
rates and probabilities. It is not an unresolved route to explosion
while the bounds of Theorem~\ref{thm:C} are retained.


\section{Why B survives bounded throughput:
specific turnover and shrinking inventory}

Consider the continuum cascade
\begin{align}
\Phi_n&=\kappa_ne_n^m,\qquad m>1, \label{eq:flux}\\
\dot e_0&=J-\Phi_0,\\
\dot e_n&=\eta\Phi_{n-1}-\Phi_n,\qquad n\ge1,
\label{eq:cascade}
\end{align}
where $J>0$, $0<\eta\le1$, and $e_n$ is the resource inventory
assigned to level $n$.

For $e_n>0$, define specific turnover and inventory residence time:
\begin{equation}
\rho_n=\frac{\Phi_n}{e_n}
=\kappa_ne_n^{m-1},
\qquad
\tau_n^{\rm res}=\frac1{\rho_n}.
\label{eq:rho}
\end{equation}

The residence-time symbol is distinguished from the actual
retained-transition duration $\tau_n$. Identifying the two requires
a model of what constitutes completion of the transformation.

\begin{theorem}[Bounded throughput and unbounded specific turnover]
\label{thm:B}
The following statements hold.

\begin{enumerate}
\item If $0<\rho_n\le\rho_{\max}<\infty$, then
\[
\sum_n\tau_n^{\rm res}=\infty.
\]
The same conclusion holds for actual durations satisfying
$\tau_n\ge a/\rho_n$ with a constant $a>0$.

\item For $\kappa_n=\kappa_0b^n$, with $\kappa_0,b>0$,
the forced stationary profile is
\[
\Phi_n^*=J\eta^n,
\qquad
e_n^*=
\left(\frac{J}{\kappa_0}\right)^{1/m}
\left(\frac{\eta}{b}\right)^{n/m}.
\]
Its specific turnover is
\begin{equation}
\rho_n^*
=
\kappa_0^{1/m}J^{(m-1)/m}
\left(b\eta^{m-1}\right)^{n/m}.
\label{eq:rho-star}
\end{equation}
The stationary residence-time series converges exactly when
\[
b\eta^{m-1}>1.
\]

\item If $0<\eta<1$ and $b\eta^{m-1}>1$, both
$\sum_ne_n^*$ and $\sum_n\Phi_n^*$ are finite, although
$\rho_n^*\to\infty$.
\end{enumerate}
\end{theorem}

\begin{proof}
For the first statement,
$\tau_n^{\rm res}\ge1/\rho_{\max}$.

At stationarity, $\Phi_0^*=J$ and
$\Phi_n^*=\eta\Phi_{n-1}^*$, giving $\Phi_n^*=J\eta^n$.
Substituting this into \eqref{eq:flux} gives the inventory formula.
Dividing throughput by inventory gives \eqref{eq:rho-star}.
The residence times form a geometric series with ratio
$(b\eta^{m-1})^{-1/m}$.

Finally, for $0<\eta<1$,
\[
\sum_n\Phi_n^*=\frac{J}{1-\eta}<\infty.
\]
The supercritical specific-turnover condition implies
$b>\eta^{1-m}>1>\eta$, so
$(\eta/b)^{1/m}<1$ and the inventory series also converges.
\end{proof}

Thus
\begin{equation}
\boxed{
\text{bounded aggregate throughput does not imply
bounded specific turnover}.
}
\label{eq:throughput-distinction}
\end{equation}

The distinction is a shrinking denominator. The rate
$\Phi_n/e_n$ can increase while the physical amount processed
at each level decreases.

For example, $m=2$, $\eta=1/2$, and $b=4$ give finite total
stationary inventory and aggregate throughput, while
$\rho_n^*$ grows proportionally to $(\sqrt2)^n$.
If $e_n$ is an energy inventory and $1-\eta$ is the dissipated
fraction per handoff, stationary dissipation is
\[
(1-\eta)\sum_n\Phi_n^*=J.
\]
This accounts for the resource transfers in the cascade.
It does not account automatically for the apparatus needed to
realize its increasingly large kinetic coefficients.

A power bound likewise need not control the ratio. If processing
requires work $w_0>0$ per unit throughput, then
\[
w_0\Phi_n\le P_{\max}
\quad\Longrightarrow\quad
\rho_n\le\frac{P_{\max}}{w_0e_n},
\]
which is not uniform when $e_n\to0$.

Conversely, a model that establishes
$\Phi_n\le k_{\max}C_n$ and
$C_n/e_n\le\chi_{\max}$ obtains
$\rho_n\le k_{\max}\chi_{\max}$.
Here $C_n$ must measure physical processing capacity in compatible
units. A global bound on $C_n$ alone does not bound its ratio to
a vanishing inventory.

Theorem~\ref{thm:B} is a stationary model calculation. It neither
proves that a transient solution reaches infinite retained depth
in finite time nor derives $\kappa_n$ from retained architecture.
An infinite stationary profile already contains nonzero inventory
at every level.

The physical question is therefore:
\begin{quote}
What permits successive retained transformations to occur at
increasing specific rate while the amount being transformed shrinks?
\end{quote}
The simple processor-count argument that closes C does not answer
this continuum question. Neither does prescribing $\kappa_n$.


\section{The finite-action theorem}

Let the physical states belong to a metric space
$(\mathsf X,d_{\rm phys})$, and write
\[
d_{\rm phys,n}=d_{\rm phys}(x_n,x_{n+1}).
\]
Suppose each transition obeys a speed--distance inequality
\begin{equation}
\varepsilon_n\tau_n
\ge c_n d_{\rm phys,n}^2.
\label{eq:action-law}
\end{equation}

The physical meaning of this inequality must be justified for the
chosen class of dynamics. Its costs, metric, and coefficients must
be specified together. The theorem below is an exact mathematical
consequence of the inequality; its use as a physical no-go theorem
is conditional on that justification.

\begin{theorem}[Finite-action bound]
\label{thm:action}
Suppose \eqref{eq:action-law} holds with
$c_n\ge c_*>0$, and
\[
\mathcal E:=\sum_n\varepsilon_n<\infty,
\qquad
\mathcal T:=\sum_n\tau_n<\infty.
\]
Then
\begin{equation}
\boxed{
\sum_n d_{\rm phys,n}
\le
\frac1{\sqrt{c_*}}\sqrt{\mathcal E\mathcal T}
<\infty.
}
\label{eq:finite-length}
\end{equation}
\end{theorem}

\begin{proof}
For every finite $k$,
\[
\begin{aligned}
\sum_{n<k}d_{\rm phys,n}
&\le
\frac1{\sqrt{c_*}}
\sum_{n<k}\sqrt{\varepsilon_n\tau_n}\\
&\le
\frac1{\sqrt{c_*}}
\sqrt{
\left(\sum_{n<k}\varepsilon_n\right)
\left(\sum_{n<k}\tau_n\right)
}\\
&\le
\frac1{\sqrt{c_*}}\sqrt{\mathcal E\mathcal T}.
\end{aligned}
\]
The second inequality is Cauchy--Schwarz. Taking the limit of the
nondecreasing partial sums proves the result.
\end{proof}

Consequently,
\[
\text{finite charged energy}
+\text{finite time}
+\text{uniform speed--distance control}
\]
implies finite accumulated physical distance.

No assumption of monotonic organizational improvement is used.
The theorem applies equally to productive changes, reversals,
oscillations, or unsuccessful transformations, provided the stated
physical inequalities apply to the counted transitions.


\section{The fixed-resolution transition-count corollary}

For $\delta>0$, let
\[
K_{\delta}^{\rm phys}(T)
=
\#\left\{
n:t_{n+1}\le T,\;
d_{\rm phys,n}\ge\delta
\right\}.
\]

\begin{corollary}[Fixed physical resolution]
\label{cor:physical}
Under the uniform-coefficient assumptions of
Theorem~\ref{thm:action} and the budget \eqref{eq:budget},
\begin{equation}
\boxed{
K_{\delta}^{\rm phys}(T)
\le
\frac1\delta
\sqrt{\frac{B(T)T}{c_*}}.
}
\label{eq:physical-count}
\end{equation}
\end{corollary}

\begin{proof}
For any finite collection of qualifying completed transitions,
their distances sum to at least their number times $\delta$.
Their costs sum to at most $B(T)$, and their disjoint durations
sum to at most $T$. The finite-prefix argument in
Theorem~\ref{thm:action} therefore gives
\[
K\delta\le\sqrt{B(T)T/c_*}.
\]
The same bound for every finite collection excludes an infinite
qualifying collection.
\end{proof}

For a general organizational conclusion, introduce an operational distance
$d_{\rm op}$ and suppose that the relevant physical metric satisfies
\begin{equation}
d_{\rm phys}(x,y)\ge\alpha_*d_{\rm op}(x,y),
\qquad \alpha_*>0.
\label{eq:metric-comparison}
\end{equation}

\begin{corollary}[Fixed operational resolution]
\label{cor:operational}
Under the assumptions of Corollary~\ref{cor:physical} and
\eqref{eq:metric-comparison},
\begin{equation}
\boxed{
K_{\delta}^{\rm op}(T)
\le
\frac1{\alpha_*\delta}
\sqrt{\frac{B(T)T}{c_*}},
}
\label{eq:operational-count}
\end{equation}
where $K_{\delta}^{\rm op}(T)$ counts completed \emph{consecutive}
transitions with $d_{\rm op}(x_n,x_{n+1})\ge\delta$. This is a path-length
or transition-count quantity; Section~\ref{sec:packing-depth} shows why it
must not be identified directly with retained organizational depth.
\end{corollary}

\begin{proof}
Every counted transition has
$d_{\rm phys,n}\ge\alpha_*\delta$.
Apply Corollary~\ref{cor:physical}.
\end{proof}

Equation~\eqref{eq:metric-comparison} is an additional assumption in the
general theorem, not in every model. Section~\ref{sec:metrics} proves the
comparison with sharp constant one for finite-state Markov dynamics when the
physical distance is total variation and operational readout is represented by
a stochastic channel. For transport metrics the comparison requires additional
readout regularity and is not automatic.

The counting argument applies to any disjoint sequential partition
for which the cost budget and speed--distance law hold. In the
uniform-coefficient case, grouping adjacent transitions preserves
the corresponding bound by the triangle inequality and
Cauchy--Schwarz. Arbitrary subdivision therefore does not create
an unlimited number of fixed-resolution changes.

The converse interpretation requires care. A large number of
fixed-resolution consecutive changes need not correspond to a large number
of distinct retained organizations: a trajectory may revisit the same
states. The transition-count corollary is therefore retained as a sharp
dynamical bound, but retained depth is defined separately below.

\section{Variable kinetic coefficients and the remaining boundary}

The uniform lower bound on $c_n$ is not needed for a weighted result.

\begin{corollary}[Weighted finite-action bound]
\label{cor:weighted}
If \eqref{eq:action-law} holds with $c_n\ge0$ and
$\mathcal E,\mathcal T<\infty$, then
\begin{equation}
\boxed{
\sum_n\sqrt{c_n}\,d_{\rm phys,n}
\le\sqrt{\mathcal E\mathcal T}<\infty.
}
\label{eq:weighted-length}
\end{equation}
\end{corollary}

\begin{proof}
For every finite $k$,
\[
\sum_{n<k}\sqrt{c_n}\,d_{\rm phys,n}
\le
\sum_{n<k}\sqrt{\varepsilon_n\tau_n}
\le
\sqrt{\mathcal E\mathcal T}.
\]
Take the limit.
\end{proof}

If $d_{\rm phys,n}\ge\delta>0$ for all counted transitions, then
\begin{equation}
\sum_n\sqrt{c_n}
\le
\frac1\delta\sqrt{\mathcal E\mathcal T}<\infty,
\label{eq:coefficient-sum}
\end{equation}
and therefore $c_n\to0$.

Conversely, if $c_n\ge c_*>0$, then
$\sum_nd_{\rm phys,n}<\infty$ and hence
$d_{\rm phys,n}\to0$.

These two implications do not establish the unrestricted disjunction
``$c_n\to0$ or $d_{\rm phys,n}\to0$.'' The factors can alternate.
For normalized scalar sequences, take
\[
(c_n,d_n)=
\begin{cases}
(1,2^{-n}), & n\ \text{even},\\
(4^{-n},1), & n\ \text{odd}.
\end{cases}
\]
Then $\sqrt{c_n}d_n=2^{-n}$ is summable, while neither factor
converges to zero. This is an algebraic counterexample to an
overstrong limiting statement, not a physical escape mechanism.

The exact general condition is \eqref{eq:weighted-length}.
In particular, for any fixed positive $c_0,d_0$, only finitely many
transitions can simultaneously satisfy
$c_n\ge c_0$ and $d_{\rm phys,n}\ge d_0$.

The interpretation of coefficient collapse is model-dependent.
It can correspond, in an appropriate formulation, to increasing
mobility or dynamical activity, decreasing friction, changing control
constraints, or the failure of the assumed speed-limit description.
These interpretations are not interchangeable.

Nor is an arbitrarily weakened bound evidence of physical
coefficient collapse. If a model admits a valid uniform positive
coefficient, replacing it with smaller numbers does not remove the
no-go result. A proposed escape must explain why no applicable
uniform control survives under a fixed physical metric and
consistent accounting.


\section{Sharpness example:
finite-time infinite refinement of a finite trajectory}

\begin{example}[Bounded-power refinement]
\label{ex:refinement}
Consider the deterministic overdamped model
\[
\zeta\dot x=F,
\qquad
\zeta>0,\quad F>0,
\]
on an interval containing $[0,L]$. Its speed is
$v=F/\zeta$ and its dissipative power is
\[
P=Fv=\zeta v^2.
\]

Define refinement labels
\[
x_n=L(1-2^{-n}).
\]
The dynamics contains no level-dependent acceleration coefficient.
The physical increments, crossing durations, and dissipated costs are
\[
d_n=L2^{-(n+1)},
\qquad
\tau_n=\frac{d_n}{v},
\qquad
\varepsilon_n=Fd_n.
\]
Consequently,
\[
\sum_n\tau_n=\frac Lv,
\qquad
\sum_n\varepsilon_n=FL,
\qquad
\sum_nd_n=L,
\]
and
\[
\varepsilon_n\tau_n=\zeta d_n^2.
\]
The finite-action length bound is saturated.
\end{example}

The example produces infinitely many labelled transitions that are
successively faster and cheaper at bounded power. It therefore
shows why finite time and power cannot exclude infinite indexing.

Nevertheless,
\[
x_n\to L,
\]
and the physical trajectory extends regularly through $x=L$.
No state variable or physical velocity diverges.

The choice of refinement labels is deliberate: an ordinary finite
transformation has been partitioned into infinitely many increasingly
small changes. Treating each label as a separately recorded event
would add a controller or record-keeping problem not contained in
the example.

This is a continuum sharpness example, not a construction of
infinitely many reliable organizational innovations. It shows that
the fixed-resolution condition is essential, while making no claim
that arbitrarily fine labels remain distinguishable under bounded
measurement resources.


\section{Operational distinguishability and measurement resolution}

An operational distinction must specify the allowed observation.
For a concrete classical readout, let retained states $x,y$
produce probability distributions $p_x,p_y$. With equal priors,
the optimal single-shot discrimination error is
\begin{equation}
P_{\rm err}
=
\frac{1-d_{\rm TV}(p_x,p_y)}{2},
\qquad
d_{\rm TV}(p,q)=\frac12\|p-q\|_1.
\label{eq:discrimination}
\end{equation}
A requirement $P_{\rm err}\le\epsilon<1/2$ therefore implies
\[
d_{\rm TV}(p_x,p_y)\ge1-2\epsilon>0.
\]

More generally, fix a class $\mathcal M$ of admissible measurement
channels and define
\[
d_{\rm op}(x,y)
=
\sup_{M\in\mathcal M}
d_{\rm TV}(p_x^M,p_y^M).
\]
This may be a pseudometric: physically different states can be
indistinguishable to every permitted measurement. It becomes a
metric after identifying operationally indistinguishable states.

The class $\mathcal M$ must specify the measurement standard. In
particular, observation time, apparatus, precision, number of repetitions,
and tolerated error cannot be allowed to grow without accounting. Repeated
independent samples can amplify arbitrarily small single-shot differences;
therefore fixed operational resolution means fixed measurement resources
as well as a fixed discrimination threshold.

\begin{definition}[Fixed-resolution transition activity]
For $\delta>0$, the consecutive fixed-resolution count up to time $T$ is
\[
K_\delta^{\rm consec}(T)
=
\#\{n:t_{n+1}\le T,\ d_{\rm op}(x_n,x_{n+1})\ge\delta\}.
\]
This measures resolved path activity. It does not require a counted state to
be new relative to states visited earlier.
\end{definition}

\begin{definition}[Retained packing depth]
Let $V(T)=\{x_n:t_n\le T\}$ be the visited retained states up to time $T$.
For a fixed operational resolution $\delta>0$, define the packing depth as
the supremum over finite mutually separated subfamilies
\begin{equation}
\boxed{
D_\delta(T)
=
\sup\left\{
|Y|:Y\subseteq V(T),\;Y\text{ finite},\;
d_{\rm op}(x,y)\ge\delta
\text{ for all distinct }x,y\in Y
\right\}.
}
\label{eq:packing-depth}
\end{equation}
Thus revisiting an old state can add transition activity without increasing
retained depth.
\end{definition}

\begin{definition}[Resolution-robust organizational escape]
A retained process exhibits resolution-robust organizational escape if there
exist $\delta>0$ and $T_*<\infty$ such that $D_\delta(T)$ is unbounded as
$T\uparrow T_*$ under one fixed admissible measurement standard.
\end{definition}

\begin{definition}[Vanishing-resolution infinite indexing]
An infinitely indexed retained-transition sequence has vanishing-resolution
indexing if
\[
d_{\rm op}(x_n,x_{n+1})\to0.
\]
\end{definition}

Vanishing-resolution indexing may describe refinement toward a finite state,
as in Example~\ref{ex:refinement}. A different pathology is recurrence:
a two-state shuttle can maintain
$d_{\rm op}(x_n,x_{n+1})=\delta$ forever while
$D_\delta(T)=2$. The packing-depth definition excludes both failures for
different reasons.

Operational distinguishability does not require every admissible
test to be performed. If measurement, verification, or record
updates are actually used to generate or retain the sequence,
however, their costs and timing belong inside the physical model.

\section{Transition count versus retained packing depth}
\label{sec:packing-depth}

The distinction between $K_\delta^{\rm consec}$ and $D_\delta$ is not
cosmetic. A two-state shuttle can make the consecutive fixed-resolution count
arbitrarily large while visiting only two operationally distinct states.
Transition count is resolved path activity; retained depth is the size of the
mutually distinguishable repertoire actually visited.

The finite-state activity--entropy bound is nevertheless the right dynamical
object. Its tightness can be characterized exactly.

\subsection*{The transition speed limit is asymptotically sharp}

Consider a driven edge with forward and reverse probability fluxes
\[
a=\frac{A(1+\epsilon)}2,
\qquad
b=\frac{A(1-\epsilon)}2,
\qquad
0<\epsilon<1.
\]
The net probability flux is $J=\epsilon A$ and the entropy-production rate is
\[
\sigma
=
(a-b)\log\frac ab
=
2\epsilon A\,\operatorname{artanh}\epsilon.
\]
Over a constant-flux interval of duration $\tau$,
\[
d_{\rm TV}=\epsilon A\tau,
\qquad
\mathcal N=A\tau,
\qquad
\Sigma=
2\epsilon A\tau\,\operatorname{artanh}\epsilon.
\]
Hence
\begin{equation}
\boxed{
r_{\rm SFS}(\epsilon)
=
\frac{d_{\rm TV}}{\sqrt{\mathcal N\Sigma/2}}
=
\sqrt{\frac{\epsilon}{\operatorname{artanh}\epsilon}}
<1
\quad(\epsilon>0),
\qquad
r_{\rm SFS}(\epsilon)\to1
\quad(\epsilon\to0).
}
\label{eq:sfs-tightness}
\end{equation}
The saturation is asymptotic, not exact. At fixed resolved displacement
$d_{\rm TV}=\delta$,
\[
\mathcal N=\frac{\delta}{\epsilon}\to\infty,
\qquad
\Sigma=2\delta\,\operatorname{artanh}\epsilon\to0.
\]
Thus entropy production per resolved distinction can approach zero only by
paying diverging dynamical activity.

The sharper form proposed by Zhang \cite{Zhang2018Comment} has the same
saturating construction. Counting each unordered edge once, define
\[
B=
\left(\sqrt a+\sqrt b\right)^2
=
A\left(1+\sqrt{1-\epsilon^2}\right).
\]
The corresponding achieved-to-bound ratio is
\begin{equation}
\boxed{
r_Z(\epsilon)
=
\sqrt{
\frac{2\epsilon}
{\operatorname{artanh}\epsilon\,
\left(1+\sqrt{1-\epsilon^2}\right)}
}
\longrightarrow1
\quad(\epsilon\to0).
}
\label{eq:zhang-tightness}
\end{equation}
Moreover $B/(2A)\to1$ in this same limit, so the SFS and Zhang forms coincide
where saturation occurs.

\subsection*{Saturation is independent of retained depth}

A shuttle between two distributions can repeat the tight transition
indefinitely while retaining packing depth two. Sharpness is not peculiar to
that depthless construction. For an $m$-state simplex and $0<s<1$, let
\[
q=\frac{1-s}{m},
\qquad
p^{(j)}=q\mathbf 1+s e_j.
\]
For $i\ne j$,
\[
d_{\rm TV}(p^{(i)},p^{(j)})=s.
\]
Driving exactly $s$ probability mass from one peak to the next along the same
constant-asymmetry edge gives the same thermodynamic tightness ratio. Every
coordinate remains at least $q$, so the required rates stay finite, bounded by
$a/q=a m/(1-s)$.

At measurement resolution $\delta\le s$, each move counts once. The
achieved-to-counting-bound ratio is therefore
\[
\frac{\delta}{s}\,r_{\rm SFS}(\epsilon).
\]
The earlier vertex-ring value $0.50$ at $\delta=0.5$ was simply the
$s\to1$ limit of $\delta/s$, while that same limit also drives the maximum
required rate to infinity. Choosing the efficient step $s=\delta$ removes both
artifacts: finite-depth rings and the depth-two shuttle lie on the same
tightness curve.

\subsection*{Two independent obstructions to retained depth}

The previous section defined $D_\delta(T)$ as the supremum of the cardinalities
of finite pairwise-$\delta$-separated subsets of the visited retained
repertoire. Let $\mathcal P_\delta(\mathcal X_T)$ denote the
$\delta$-packing number of the accessible operational state space.

\begin{theorem}[Two-factor retained-depth bound]
\label{thm:twofactor-depth}
For finite-state Markov jump dynamics under the operational speed-limit
hypotheses of Theorem~\ref{thm:markov-operational},
\begin{equation}
\boxed{
D_\delta(T)
\le
\min\left\{
1+\frac1\delta
\sqrt{\frac{\Sigma_*\mathcal N_*}{2}},
\;
\mathcal P_\delta(\mathcal X_T)
\right\}.
}
\label{eq:combined-depth}
\end{equation}
\end{theorem}

\begin{proof}
Choose any finite pairwise-$\delta$-separated family of $D$ visited retained
states and order their first selected visit times chronologically. The
$D-1$ intervals between consecutive selected representatives have pairwise
disjoint interiors, and every such pair has operational distance at least
$\delta$. Applying the data-processing comparison and the integrated
activity--entropy speed limit to each interval gives
\[
\delta
\le
\sqrt{\frac{\mathcal N_k\Sigma_k}{2}}.
\]
Summing over the selected intervals and applying Cauchy--Schwarz,
\[
(D-1)\delta
\le
\frac1{\sqrt2}
\left(\sum_k\mathcal N_k\right)^{1/2}
\left(\sum_k\Sigma_k\right)^{1/2}
\le
\sqrt{\frac{\mathcal N_*\Sigma_*}{2}},
\]
because activity and entropy production are non-negative and the selected
intervals are disjoint. This gives the thermodynamic factor. The geometric
factor follows directly because the selected family is a
$\delta$-separated subset of $\mathcal X_T$. Since the thermodynamic bound
holds for every finite selected family, it bounds the packing depth itself.
\end{proof}

The finite-block mathematical implication is independently machine-checked in
\texttt{PackingDepth.lean} as
\texttt{packing\_depth\_thermo} and
\texttt{packing\_depth\_two\_factor}. The speed-limit premise and the
state-space-specific packing theorem remain external assumptions/results, not
facts hidden inside the formalization.

The two factors answer different questions. The thermodynamic factor asks
whether the trajectory has enough activity--entropy budget to traverse enough
resolved distance. The geometric factor asks whether the operational state
space contains enough mutually distinguishable places to visit.

\subsection*{A tighter finite-simplex packing bound}

For probability distributions on a fixed finite state space $|S|=m$, total
variation makes the simplex $\Delta_{m-1}$ compact, so
$\mathcal P_\delta(\Delta_{m-1})$ is finite for every $\delta>0$.

A useful explicit bound comes from largest-remainder rounding to the lattice
\[
L_k
=
\left\{
\frac{z}{k}:
z\in\mathbb Z_{\ge0}^{m},
\ \sum_i z_i=k
\right\}.
\]
For every $p\in\Delta_{m-1}$ there exists a rounded point
$\pi_k(p)\in L_k$ with
\[
d_{\rm TV}(p,\pi_k(p))
\le
\frac{m}{4k}.
\]
Choose
\[
k=
\left\lfloor\frac{m}{2\delta}\right\rfloor+1.
\]
If two $\delta$-separated simplex points rounded to the same lattice point,
the triangle inequality would give
$d_{\rm TV}(p,q)\le m/(2k)<\delta$, a contradiction. Therefore rounding is
injective on every $\delta$-packing, and
\begin{equation}
\boxed{
\mathcal P_\delta(\Delta_{m-1})
\le
\binom{k+m-1}{m-1},
\qquad
k=
\left\lfloor\frac{m}{2\delta}\right\rfloor+1.
}
\label{eq:simplex-packing}
\end{equation}
For $m=2$, the exact value is
\[
\mathcal P_\delta(\Delta_1)
=
\left\lfloor\frac1\delta\right\rfloor+1.
\]
The lattice bound is substantially tighter than the coarse volumetric
estimate used in the preceding draft.

\subsection*{Both factors can bind}

Neither factor in \eqref{eq:combined-depth} is decorative. For an interior
ring with $m=12$, $\delta=0.5$, and $\epsilon=0.01$, the achieved packing depth
is
\[
D_\delta=K+1
\]
for $K<m$ transitions and agrees with the thermodynamic factor to five
figures. Once all twelve configurations have been visited, $D_\delta$ remains
at twelve while the thermodynamic factor continues to increase. The geometric
factor then becomes the active constraint.

Figure~\ref{fig:twofactor-tightness} displays both regimes. Panel (a) shows
that a depth-two shuttle and a depth-six interior ring have the same
asymptotic tightness curve. Panel (b) shows the crossover from a
thermodynamic-resource ceiling to the fixed repertoire ceiling.

\begin{figure}[t]
\centering
\includegraphics[width=\textwidth]{verification/organizational-depth/tightness/figures/tightness.pdf}
\caption{Sharpness of the transition bound and the two-factor retained-depth
boundary. \textbf{(a)} The SFS and Zhang achieved-to-bound ratios as functions
of drive asymmetry, with measured points from a depth-two shuttle and a
depth-six finite-rate interior ring. Saturation is asymptotic and independent
of retained depth. \textbf{(b)} Packing depth for a twelve-configuration
interior ring as activity budget increases. The thermodynamic factor binds
until the finite repertoire is exhausted; the geometric factor binds
afterward.}
\label{fig:twofactor-tightness}
\end{figure}

\begin{corollary}[What unbounded fixed-resolution depth requires]
\label{cor:unbounded-packing-depth}
At fixed $\delta>0$, unbounded retained packing depth requires both
\[
\mathcal P_\delta(\mathcal X)=\infty
\]
and an unbounded activity--entropy product along the history. Either finite
operational packing capacity or a uniform finite ceiling on
$\mathcal N_*\Sigma_*$ forbids unbounded depth.
\end{corollary}

The geometric requirement is more general than growth of a finite-state
alphabet. A fixed countably infinite state space, an infinite-dimensional
state description, a non-compact memory coordinate, or a physical record that
grows without bound can all provide infinite fixed-resolution packing
capacity. What matters is not literally increasing $|S|$, but the physical
creation and maintenance of an operational repertoire that is not totally
bounded at the declared resolution.

\section{Physical versus organizational metrics}
\label{sec:metrics}

The physical metric in a speed limit and the operational metric
of organizational distinction need not coincide. A uniform comparison
$d_{\rm phys}\ge\alpha_*d_{\rm op}$ is therefore a genuine condition of
the general theorem. It can, however, be a theorem rather than an assumption
in a concrete model.

\subsection*{Finite-state Markov dynamics: the bridge is exact}

Let $S$ be a finite state space and let $p,q$ be probability
distributions on $S$. A measurement channel is a stochastic kernel
$\Phi(o\mid i)$ from states $i$ to outcomes $o$. Its readout distribution
is $\Phi p$. Total variation satisfies the data-processing inequality
\begin{equation}
d_{\rm TV}(\Phi p,\Phi q)\le d_{\rm TV}(p,q).
\label{eq:tv-data-processing}
\end{equation}
Indeed,
\[
\sum_o\left|\sum_i\Phi(o\mid i)(p_i-q_i)\right|
\le
\sum_i|p_i-q_i|\sum_o\Phi(o\mid i)
=
\sum_i|p_i-q_i|.
\]
Taking the supremum over any fixed admissible measurement class gives
\begin{equation}
\boxed{
d_{\rm op}(p,q)\le d_{\rm TV}(p,q).
}
\label{eq:tv-operational}
\end{equation}
The universal constant one is sharp over unrestricted channels, because the
identity channel attains equality.

Now consider a continuous-time Markov jump process with rates
$W_{ij}(t)$ for the transition $j\to i$ and distribution $p(t)$, under the
standard entropy-production conditions of the Shiraishi--Funo--Saito bound
(for example, a thermodynamic interpretation with local detailed balance). Define the
directed probability flux
\[
F_{ij}(t)=W_{ij}(t)p_j(t),
\]
the dynamical activity
\[
A(t)=\sum_{i<j}\bigl(F_{ij}(t)+F_{ji}(t)\bigr),
\]
and the dimensionless entropy-production rate
\[
\sigma(t)
=
\sum_{i<j}
\bigl(F_{ij}-F_{ji}\bigr)
\ln\frac{F_{ij}}{F_{ji}},
\]
with the usual limiting interpretation when a reverse flux vanishes.
For an interval $I$, write
\[
\mathcal N(I)=\int_I A(t)\,dt,
\qquad
\Sigma(I)=\int_I\sigma(t)\,dt.
\]
Here $\mathcal N(I)$ is the expected number of jumps on the interval.

The Shiraishi--Funo--Saito speed limit gives
\begin{equation}
d_{\rm TV}\bigl(p(t),p(t')\bigr)^2
\le
\frac12\,\mathcal N(I)\Sigma(I).
\label{eq:sfs-tv}
\end{equation}
Combining \eqref{eq:tv-operational} and \eqref{eq:sfs-tv} yields
\begin{equation}
\boxed{
2d_{\rm op}\bigl(p(t),p(t')\bigr)^2
\le
\mathcal N(I)\Sigma(I).
}
\label{eq:operational-action}
\end{equation}

\begin{theorem}[Operational activity--entropy bound]
\label{thm:markov-operational}
Let $I_n=[t_n,t_{n+1}]$ be disjoint sequential transition intervals of a
finite-state Markov jump process satisfying the conditions of
\eqref{eq:sfs-tv}, and let
$d_{{\rm op},n}=d_{\rm op}(p(t_n),p(t_{n+1}))$.
If
\[
\sum_n\Sigma(I_n)\le\Sigma_*<\infty,
\qquad
\sum_n\mathcal N(I_n)\le\mathcal N_*<\infty,
\]
then
\begin{equation}
\boxed{
\sum_n d_{{\rm op},n}
\le
\sqrt{\frac{\Sigma_*\mathcal N_*}{2}}.
}
\label{eq:markov-length}
\end{equation}
Consequently, for every $\delta>0$,
\begin{equation}
\boxed{
K_\delta^{\rm op}
\le
\frac1\delta
\sqrt{\frac{\Sigma_*\mathcal N_*}{2}}.
}
\label{eq:markov-count}
\end{equation}
\end{theorem}

\begin{proof}
Equation~\eqref{eq:operational-action} gives
$d_{{\rm op},n}\le
2^{-1/2}\sqrt{\Sigma_n\mathcal N_n}$.
Summing over a finite prefix and applying Cauchy--Schwarz gives
\[
\sum_n d_{{\rm op},n}
\le
\frac1{\sqrt2}
\sqrt{\sum_n\Sigma_n}
\sqrt{\sum_n\mathcal N_n},
\]
and monotone convergence gives \eqref{eq:markov-length}.
Every transition counted by $K_\delta^{\rm op}$ contributes at least
$\delta$ to the left-hand side, giving \eqref{eq:markov-count}.
\end{proof}

\begin{corollary}[Thermodynamic retained-depth bound]
\label{cor:markov-depth}
Under the hypotheses of Theorem~\ref{thm:markov-operational}, let
$y_0,\ldots,y_{D-1}$ be retained states encountered in chronological order,
with
\[
d_{\rm op}(y_i,y_j)\ge\delta>0
\qquad(i\ne j).
\]
Then
\begin{equation}
\boxed{
D
\le
1+\frac1\delta
\sqrt{\frac{\Sigma_*\mathcal N_*}{2}}.
}
\label{eq:markov-depth}
\end{equation}
Consequently the retained packing depth obeys the same bound.
\end{corollary}

\begin{proof}
There are $D-1$ disjoint chronological intervals between successive selected
representatives. Pairwise separation implies that each such interval has
endpoint operational distance at least $\delta$. Apply
Theorem~\ref{thm:markov-operational} to those aggregated intervals:
\[
(D-1)\delta
\le
\sqrt{\Sigma_*\mathcal N_*/2}.
\]
Rearranging gives \eqref{eq:markov-depth}. The corresponding finite-prefix
statement is machine-checked in
\texttt{OperationalBridge.lean} as
\texttt{pairwise\_depth\_resource\_bound}.
\end{proof}

This result has a useful interpretation in the notation of the general
finite-action theorem. If $I_n$ has duration $\tau_n$ and average activity
$\bar A_n=\mathcal N_n/\tau_n$, then \eqref{eq:sfs-tv} implies
\[
\Sigma_n\tau_n
\ge
\frac{2}{\bar A_n}d_{{\rm TV},n}^2
\ge
\frac{2}{\bar A_n}d_{{\rm op},n}^2.
\]
Under an isothermal interpretation in which the charged dissipative cost is
$\varepsilon_n=k_BT\Sigma_n$,
\begin{equation}
\varepsilon_n\tau_n
\ge
\frac{2k_BT}{\bar A_n}d_{{\rm op},n}^2.
\label{eq:activity-coefficient}
\end{equation}
Thus, in this model class,
\[
c_n=\frac{2k_BT}{\bar A_n}.
\]
Loss of a uniform kinetic coefficient corresponds concretely to unbounded
average dynamical activity. More recent thermodynamic speed-limit work likewise
identifies activity or mobility as the kinetic resource controlling attainable
speed \cite{VuSaito2023,Nagayama2025}.

\subsection*{Transport metrics: no unconditional TV comparison}

The corresponding statement is not automatic for Wasserstein distance. There
is no universal $\alpha>0$ such that
\[
W_2(\mu,\nu)\ge\alpha d_{\rm TV}(\mu,\nu)
\]
for arbitrary probability measures: for
$\mu=\delta_0$ and $\nu=\delta_\epsilon$,
$W_2(\mu,\nu)=\epsilon\to0$ while
$d_{\rm TV}(\mu,\nu)=1$.

This does not make transport-based operational bounds impossible; it means
additional readout regularity is required. Suppose, for every admissible
measurement $M$ and output event $B$, the conditional readout probability
\[
g_{M,B}(x)=\Pr_M(B\mid x)
\]
has Lipschitz constant at most $L_{\rm readout}$ uniformly in $M,B$. Then a
coupling argument gives
\[
d_{\rm op}(\mu,\nu)
\le L_{\rm readout}W_1(\mu,\nu)
\le L_{\rm readout}W_2(\mu,\nu).
\]
Hence a $W_2$ speed limit does transfer under an explicit bounded-sensitivity
measurement assumption. The distinction is between an unconditional global
comparison, which fails, and a conditional readout comparison, which can hold.

\subsection*{Metric decoupling outside the finite-state TV class}

On transitions with $d_{\rm op,n}>0$, define
\[
\alpha_n=
\frac{d_{\rm phys,n}}{d_{\rm op,n}}.
\]
If operationally significant differences remain bounded below while
$\alpha_n\to0$, their physical separation collapses relative to their
operational separation. For a sequence with
$d_{\rm op,n}\ge\delta>0$, Corollary~\ref{cor:weighted} implies
\begin{equation}
\boxed{
\sum_n\sqrt{c_n}\,\alpha_n
\le
\frac1\delta\sqrt{\mathcal E\mathcal T}<\infty.
}
\label{eq:operational-weighted}
\end{equation}

Thus a uniform lower bound on $\alpha_n$ forces
$\sum_n\sqrt{c_n}<\infty$, while a uniform lower bound on $c_n$ forces
$\sum_n\alpha_n<\infty$. Without either lower bound only the joint weighted
condition follows.

Metric decoupling remains a candidate boundary outside the finite-state TV
setting. It may reflect genuine causal leverage, omitted downstream machinery,
increasing measurement resources, an expanding state description, or a metric
that loses sensitivity to the relevant physical differences. Changing
coordinates or using level-dependent normalizations cannot establish an escape:
metrics, units, and coefficients must be transformed consistently.

\subsection*{Weak convergence and stronger observables}

Finite accumulated distance makes $(x_n)$ a Cauchy sequence.
If the relevant physical metric space is complete, it converges
to a limit in that space. This concerns the retained-state
sequence; controlling every intermediate excursion requires an
appropriate trajectory-length or regularity bound.

Moreover, convergence in a weak metric does not imply boundedness
of stronger observables. Let
\[
x^{(N)}
=
(2^{-1},2^{-2},\ldots,2^{-N},0,0,\ldots).
\]
Then
\[
\|x^{(N+1)}-x^{(N)}\|_{\ell^2}=2^{-(N+1)},
\]
so the accumulated $\ell^2$ distance is finite and
$x^{(N)}$ converges in $\ell^2$. But
\[
Q(x^{(N)})
=
\sum_{j\ge1}4^j|x_j^{(N)}|^2
=N\to\infty.
\]
The limit belongs to $\ell^2$ but not to the space where this
weighted quadratic observable is finite.

Such distinctions matter for PDEs and cascades
\cite{PierreSchmitt1997,KatzPavlovic2005}.
A limit may be regular in an ambient weak space and singular
in a physically relevant stronger class. The relevant question is whether
the stronger organizational observable is operationally meaningful and
resolvable with bounded physical resources. Finite metric length alone does
not answer it.

\section{Landauer, reversible computation, and other special cases}

The finite-action theorem does not assume that every
organizational transformation erases information. Landauer's
principle supplies one conditional energetic floor; it is not
the foundation of the general no-go result.

\subsection*{Landauer as a conditional floor}

Under an appropriate isothermal erasure protocol, a net reduction
of memory entropy by $I_n$ bits gives a work or heat bound of the
form
\[
\langle W_n\rangle
\ge k_BT_n\ln2\,I_n.
\]
Its interpretation requires the memory Hamiltonian, initial
correlations, accessible side information, bath, and auxiliary
records to be accounted for
\cite{Landauer1961,ReebWolf2014}.

For example, when endpoint internal energies agree and the
controller is restored, the work and bath-heat bounds coincide.
If additionally
\[
T_n\ge T_{\min}>0,\qquad I_n\ge I_{\min}>0,
\]
a positive mean erasure-work floor follows.

Organizational change does not automatically satisfy the
erasure premise. Reversible computation may preserve
information in intermediate records and later remove those
records through uncomputation \cite{Bennett1973,Bennett1989}.
Landauer work or heat must also be distinguished from excess
dissipation, and an average erasure bound is not automatically
a fixed pathwise charge.

\subsection*{Activity-based stochastic speed limits}

Section~\ref{sec:metrics} used the exact Shiraishi--Funo--Saito
finite-state Markov bound to obtain the operational theorem
\eqref{eq:markov-count}. Its essential extra resource is dynamical
activity: entropy production constrains irreversibility, while activity
sets the kinetic time scale \cite{Shiraishi2018}.

This distinction is now supported by broader results. Discrete optimal
transport can be formulated using dynamical state mobility
\cite{VuSaito2023}, and a general activity framework generates families of
thermodynamic speed limits for Markov jump processes and deterministic
chemical reaction networks \cite{Nagayama2025}. These results reinforce the
interpretation of coefficient collapse as a kinetic-resource question, not
merely an energetic one.

\subsection*{Overdamped transport and thermodynamic length}

For an appropriate isothermal overdamped diffusion with constant
friction $\zeta>0$, probability density $p$, and probability-current
velocity $v$, the dissipative action has the form
\[
\varepsilon_{\rm diss}
=
\zeta\int_0^\tau\int p|v|^2\,dx\,dt.
\]
For admissible probability curves with finite second moments,
the transport-action inequality gives
\[
\varepsilon_{\rm diss}\tau
\ge
\zeta W_2(p_0,p_\tau)^2.
\]
This connects finite-time thermodynamics with optimal transport
\cite{Aurell2011}. As shown in Section~\ref{sec:metrics}, $W_2$ does not
globally control raw total variation; an operational transfer requires a
specified readout class or additional regularity.

Classical thermodynamic-length bounds are asymptotic or
linear-response statements \cite{Crooks2007,SivakCrooks2012}. More recent
work extends the geometry/optimal-transport connection beyond linear response
for important overdamped classes \cite{ZhongDeWeese2024}. None of these
results supplies one universal metric-speed law for arbitrary organizational
dynamics; the model and metric must remain explicit.

\subsection*{Reversible dynamics and quantum speed limits}

Logical reversibility does not by itself imply arbitrarily
fast computation. Conversely, zero thermodynamic entropy
production does not imply zero physical motion.

For example, isolated Hamiltonian evolution need not obey a
positive dissipative inequality of the particular form
\eqref{eq:action-law}. Quantum speed limits constrain evolution
using quantities such as energy above a ground state or energy
uncertainty under their respective assumptions
\cite{MargolusLevitin1998}.

Such stored dynamical resources are not automatically additive
energetic costs consumed anew at every transition. They cannot
be substituted for $\varepsilon_n$ without an accounting
argument. Bounded generator strength or energy uncertainty
can supply other metric-speed bounds, but these remain
model-specific statements.

The no-go boundary is therefore a family of conditional results,
not a claim that one dissipation formula covers all physical
implementations.

\section{Implications for Organizational Accessibility}

Organizational Accessibility separates retained architecture,
candidate exposure, the variation process, establishment, and
access to organization developed elsewhere \cite{OA8}.
This decomposition remains useful, but its singularity
interpretation must respect the bounds established here.

In the sequential stochastic model,
$\lambda_n=N_n\nu_np_n$ reveals precisely why the stated processor
and rate bounds close C. Better establishment probability cannot
grow without bound, and architectural improvements cannot restore
explosion while all hypotheses of Theorem~\ref{thm:C} remain valid.

In the continuum cascade, increasing specific turnover may
result from shrinking transformed quantities. A proposed
organizational interpretation must therefore identify the
actual retained states, their transition durations, charged
costs, physical distances, and operational distinctions.
The stationary ratio $e_n/\Phi_n$ supplies none of these
identifications automatically.

The packing-depth correction sharpens that requirement. Repeated transition
activity through a fixed repertoire is not cumulative organizational depth.
A process contributes to $D_\delta$ only by adding a retained state that
remains distinguishable from the selected repertoire at the declared
operational resolution. In this sense, retained depth is a property of the
reachable operational repertoire, not of the number of transitions executed.

Reuse, templating, catalytic machinery, recombination, and
horizontal transfer can reduce the work that remains to be
performed. They may enlarge the accessible packing repertoire, reduce the
resources needed to traverse it, or both. But neither high turnover nor a
large transition count establishes depth unless new retained distinctions
are actually added.

The ordinary cumulative claim is not threatened by this boundary. Retained
structure can improve access to later targets without generating infinitely
many targets in finite time. Indeed the corrected formulation makes the link
cleaner: cumulative emergence changes the repertoire of retained distinctions,
whereas repeated circulation through an unchanged repertoire does not.
The interventional criterion for cumulative contribution in Organizational
Accessibility remains a separate causal question from finite-time
organizational escape.

\section{Novelty boundary}

None of the following is claimed as new generic mathematics:
the pure-birth criterion, bounded-rate non-explosion,
geometric-series tests for stationary cascades,
Cauchy--Schwarz bounds, total-variation data processing,
finite-dimensional packing bounds, finite-length Cauchy convergence,
thermodynamic speed limits, optimal transport, or Zeno accumulation
\cite{Feller1971,Shiraishi2018,Aurell2011,VuSaito2023,Nagayama2025,Zhang2001}.
Nor is finite material alongside stronger-norm divergence a new
singularity mechanism \cite{KatzPavlovic2005}.

The proposed contribution is an application and clarification:
\begin{enumerate}
\item the stated physical processor assumptions already close
      the sequential stochastic accessibility route;
\item bounded aggregate throughput and bounded specific turnover
      answer different questions;
\item consecutive fixed-resolution transition count and retained
      organizational depth are different quantities;
\item retained depth can be operationalized as the packing depth of the
      visited retained repertoire at a fixed measurement resolution;
\item for finite-state Markov jump dynamics the physical-to-operational
      bridge is exact because the speed-limit metric is total variation
      and measurement channels contract total variation;
\item the activity--entropy path-length bound is asymptotically sharp, but
      sharpness alone does not diagnose depth because both recurrent and
      finite-depth trajectories can approach it;
\item retained packing depth is constrained independently by dynamical
      resources and by the packing capacity of the accessible operational
      state space;
\item outside that class, transport/readout comparison, kinetic coercivity,
      and growth of operational packing capacity remain explicit issues whose
      failure modes must be identified rather than hidden.
\end{enumerate}

No priority claim is made for the component inequalities, packing arguments,
or thermodynamic speed limits. The proposed novelty lies in assembling them
into a distinction between path activity and retained organizational depth,
and in identifying the two independent requirements for unbounded depth:
enough dynamical resources to traverse distinctions and enough operational
state-space capacity to contain them.

\section{Challenge to a dynamically organizing system}

The paper therefore ends with a constructive challenge rather than a
conjecture.

\begin{quote}
Produce a physically specified autonomous system whose retained packing depth
$D_\delta(T)$ becomes unbounded before a finite time $T_*$ at one fixed
operational resolution $\delta>0$, and provide a complete accounting of the
kinetic, energetic, control, memory, measurement, and state-space resources
that make this possible.
\end{quote}

This is stricter than asking for infinitely many resolved transitions. A
two-state shuttle already supplies those. A successful construction must add
unboundedly many mutually distinguishable retained states, or identify an
operational state variable in which such distinctions accumulate.

For the general speed--distance framework, such a construction must evade the
sufficient conditions established above by losing uniform kinetic coercivity,
losing physical-to-operational metric control, leaving the finite-cost/finite-
time regime, or moving to a state space whose fixed-resolution packing
capacity grows without bound.

For finite-state Markov jump dynamics the challenge is sharper still.
Corollary~\ref{cor:markov-depth} bounds any pairwise-separated retained
repertoire by the activity--entropy budget, while
\eqref{eq:simplex-packing} caps the repertoire geometrically for every fixed
$|S|$. Thus unbounded fixed-resolution depth requires departure from at least
one of these conditions: unbounded activity/entropy resources, failure of the
finite-state Markov description, or unbounded fixed-resolution operational packing capacity.

The last possibility turns the singularity question into a resource question
rather than resolving it for free. It need not literally be a growing finite
state count: a fixed non-totally-bounded operational space can already have
infinite packing capacity. A growing memory, record, component alphabet,
spatial support, or other operational coordinate can enlarge the packing
capacity, but the physical process that creates and maintains that capacity
must then be included in the model.

This is left deliberately open. A successful construction would be
scientifically informative precisely because it would expose both where new
distinguishable states live and what resources make them reachable and
retainable.

\section{Conclusion}

Infinite labels do not establish infinite physical change. Infinite resolved
transitions do not establish infinite retained depth. The distinction matters
because a trajectory can shuttle forever between a finite set of states while
remaining separated at fixed operational resolution on every step.

The stochastic accessibility route is closed under the stated material and
per-realization rate bounds. The continuum cascade survives bounded aggregate
throughput through a shrinking inventory and imposed increasing specific
turnover, without thereby establishing a physical transient escape.

The finite-action results remain valid as path-length and transition-count
bounds. For finite-state Markov jump dynamics the activity--entropy bound is
asymptotically sharp:
\[
\frac{d_{\rm TV}}{\sqrt{\mathcal N\Sigma/2}}
\to1
\]
for a weakly driven edge. But this sharpness is neutral about depth: both a
two-state shuttle and an interior finite-state ring can approach it.

The corrected depth quantity is the fixed-resolution packing depth
$D_\delta$, the size of the largest mutually distinguishable retained
repertoire visited by the process. It obeys two independent constraints:
\[
\boxed{
D_\delta(T)
\le
\min\left\{
1+\frac1\delta\sqrt{\frac{\Sigma_*\mathcal N_*}{2}},
\;
\mathcal P_\delta(\mathcal X_T)
\right\}.
}
\]
The first term is dynamical: enough operational path length must be paid for.
The second is geometric: enough distinguishable states must exist. For a fixed
$m$-state probability simplex,
\[
\mathcal P_\delta
\le
\left(1+\frac2\delta\right)^{m-1},
\]
so finite-state retained depth is capped regardless of how much transition
activity is spent.

The bounded-power refinement example shows why resolution is essential. The
shuttle shows why recurrence must also be excluded. The weak-versus-strong norm
example shows why finite metric length must not be mistaken for a universal
regularity theorem. The transport example shows why an operational comparison
cannot be inferred from a physical metric without specifying the readout.

A physically meaningful counterexample must therefore do more than assign an
increasing level number, prescribe increasing rates, or accumulate resolved
motion. It must create and maintain unbounded operational packing capacity and
supply the resources needed to reach new retained distinctions. That sharper
construction is left as a challenge to any proposed dynamically organizing
system.

\paragraph{Reference status.}
The supplied Organizational Accessibility v8 manuscript is the source context
for this paper. Bibliographic details and the model-specific claims used for
the stochastic speed-limit, optimal-transport, thermodynamic-length, and
activity/mobility arguments were checked against primary or publisher records
for this revision. The references are used only within the scope stated in the
text; no generic physical speed limit is inferred from a model-specific one.

\begin{thebibliography}{99}

\bibitem{OA8}
\emph{Organizational Accessibility: From Self-Maintenance to
Evolvability. A Dynamical Framework for Cumulative Emergence}.
Scientific master v8, supplied working manuscript,
12 September 2026.

\bibitem{Feller1971}
W. Feller.
\emph{An Introduction to Probability Theory and Its Applications},
Volume II, second edition.
Wiley, 1971.

\bibitem{Landauer1961}
R. Landauer.
Irreversibility and heat generation in the computing process.
\emph{IBM Journal of Research and Development}
5 (1961), 183--191.
\url{https://doi.org/10.1147/rd.53.0183}

\bibitem{Bennett1973}
C. H. Bennett.
Logical reversibility of computation.
\emph{IBM Journal of Research and Development}
17 (1973), 525--532.
\url{https://doi.org/10.1147/rd.176.0525}

\bibitem{Bennett1989}
C. H. Bennett.
Time/space trade-offs for reversible computation.
\emph{SIAM Journal on Computing}
18 (1989), 766--776.
\url{https://doi.org/10.1137/0218053}

\bibitem{ReebWolf2014}
D. Reeb and M. M. Wolf.
An improved Landauer principle with finite-size corrections.
\emph{New Journal of Physics}
16 (2014), 103011.
\url{https://doi.org/10.1088/1367-2630/16/10/103011}

\bibitem{Shiraishi2018}
N. Shiraishi, K. Funo, and K. Saito.
Speed limit for classical stochastic processes.
\emph{Physical Review Letters}
121 (2018), 070601.
\url{https://doi.org/10.1103/PhysRevLett.121.070601}

\bibitem{Aurell2011}
E. Aurell, C. Mej\'ia-Monasterio, and
P. Muratore-Ginanneschi.
Optimal protocols and optimal transport in stochastic thermodynamics.
\emph{Physical Review Letters}
106 (2011), 250601.
\url{https://doi.org/10.1103/PhysRevLett.106.250601}

\bibitem{Crooks2007}
G. E. Crooks.
Measuring thermodynamic length.
\emph{Physical Review Letters}
99 (2007), 100602.
\url{https://doi.org/10.1103/PhysRevLett.99.100602}

\bibitem{SivakCrooks2012}
D. A. Sivak and G. E. Crooks.
Thermodynamic metrics and optimal paths.
\emph{Physical Review Letters}
108 (2012), 190602.
\url{https://doi.org/10.1103/PhysRevLett.108.190602}

\bibitem{VuSaito2023}
T. Van Vu and K. Saito.
Thermodynamic unification of optimal transport: thermodynamic uncertainty
relation, minimum dissipation, and thermodynamic speed limits.
\emph{Physical Review X}
13 (2023), 011013.
\url{https://doi.org/10.1103/PhysRevX.13.011013}

\bibitem{ZhongDeWeese2024}
A. Zhong and M. R. DeWeese.
Beyond linear response: equivalence between thermodynamic geometry and
optimal transport.
\emph{Physical Review Letters}
133 (2024), 057102.
\url{https://doi.org/10.1103/PhysRevLett.133.057102}

\bibitem{Nagayama2025}
R. Nagayama, K. Yoshimura, and S. Ito.
Infinite variety of thermodynamic speed limits with general activities.
\emph{Physical Review Research}
7 (2025), 013307.
\url{https://doi.org/10.1103/PhysRevResearch.7.013307}

\bibitem{MargolusLevitin1998}
N. Margolus and L. B. Levitin.
The maximum speed of dynamical evolution.
\emph{Physica D}
120 (1998), 188--195.
\url{https://doi.org/10.1016/S0167-2789(98)00054-2}

\bibitem{Zhang2018Comment}
Y. Zhang.
Comment on \emph{Speed Limit for Classical Stochastic Processes}.
arXiv:1811.06978, 2018.
\url{https://arxiv.org/abs/1811.06978}

\bibitem{Zhang2001}
J. Zhang, K. H. Johansson, J. Lygeros, and S. Sastry.
Zeno hybrid systems.
\emph{International Journal of Robust and Nonlinear Control}
11 (2001), 435--451.
\url{https://doi.org/10.1002/rnc.592}

\bibitem{KatzPavlovic2005}
N. H. Katz and N. Pavlovi\'c.
Finite time blow-up for a dyadic model of the Euler equations.
\emph{Transactions of the American Mathematical Society}
357 (2005), 695--708.
\url{https://doi.org/10.1090/S0002-9947-04-03532-9}

\bibitem{PierreSchmitt1997}
M. Pierre and D. Schmitt.
Blowup in reaction--diffusion systems with dissipation of mass.
\emph{SIAM Journal on Mathematical Analysis}
28 (1997), 259--269.
\url{https://doi.org/10.1137/S0036141095295437}

\end{thebibliography}

% ---------------------------------------------------------------------------
% Appendix A: machine-verified mathematical core.
% Drop in after the main text, inside an \appendix environment.
% Requires amsmath, amssymb and hyperref.
% ---------------------------------------------------------------------------

\appendix
\section{Machine-verified mathematical core}
\label{app:lean}
\begingroup
\setlength{\emergencystretch}{3em}

The results of Sections~3--7 and 10 divide cleanly into two kinds of claim. One kind is
physical: that a speed--distance law of the form
$\varepsilon_n\tau_n\ge c_nd_{\rm phys,n}^2$ applies to a given class of
dynamics, that a chosen cost is legitimately charged against a cumulative
budget, and that a chosen physical metric controls operational
distinguishability. The other kind is mathematical: that the stated inequalities
entail the stated conclusions. Only the second kind can be machine-checked, and
this appendix reports that it has been.

\begin{quote}
The Lean formalization verifies the elementary mathematical implications
underlying the main no-go results, conditional on the stated hypotheses. It does
not formalize or validate the physical applicability of those hypotheses.
\end{quote}

Concretely, the mathematical core of the no-go argument was independently
formalized and machine-checked in Lean~4 with \texttt{Mathlib}. The development
contains no \texttt{sorry}, and \verb|#print axioms| reports that every result
below depends only on Lean's three standard axioms --- propositional
extensionality, quotient soundness, and choice --- so nothing rests on an
assumption added for convenience.

\subsection{What was formalized}

\paragraph{A.1 The material bound on the realization count.}
If $m_{\min}>0$ and, for every $n$, $m_{\min}N_n\le M_n$ and $M_n\le M_*$, then
\[
N_n\le\frac{M_*}{m_{\min}}\qquad\text{for every }n.
\]

\paragraph{A.2 The uniform rate bound.}
If $0\le N_n\le N_{\max}$, $0\le\nu_n\le\nu_{\max}$ and $0\le p_n\le1$, then
\[
\lambda_n=N_n\nu_np_n\le N_{\max}\nu_{\max}\qquad\text{for every }n.
\]
Only the upper bound $p_n\le1$ is used; no lower bound on $p_n$ is required.

\paragraph{A.3 The reciprocal-series obstruction (Theorem~3.1).}
If $0<\lambda_n\le\Lambda$ for every $n$, then
\[
\sum_n\frac1{\lambda_n}\text{ diverges}.
\]
This is formalized both in finite-prefix form and as
$\neg\,\mathrm{Summable}\,(n\mapsto1/\lambda_n)$. A combined declaration
derives the same conclusion directly from the material budget, material floor,
generation-rate ceiling, and $0\le p_n\le1$.

\paragraph{A.4 The weighted finite-action bound (Corollary~7.1).}
For nonnegative sequences $c_n,\varepsilon_n,\tau_n,d_n$ with
$c_nd_n^2\le\varepsilon_n\tau_n$ and summable $\varepsilon$ and $\tau$,
\[
\sum_n\sqrt{c_n}\,d_n
\le
\sqrt{\sum_n\varepsilon_n}\sqrt{\sum_n\tau_n}.
\]

\paragraph{A.5 The finite-action bound (Theorem~5.1).}
If in addition $c_n\ge c_*$ for a constant $c_*>0$, then $d$ is summable and
\[
\sum_n d_n
\le
\frac1{\sqrt{c_*}}
\sqrt{\sum_n\varepsilon_n}\sqrt{\sum_n\tau_n}.
\]
The exact variable-coefficient Lean declaration is given in the source as
\nolinkurl{finite_action_of_kinetic_floor}; the shorter
\texttt{finite\_action} is its constant-coefficient helper.

\paragraph{A.6 The fixed-resolution counting bound (Corollary~6.1).}
Under the hypotheses of A.5, for every finite set $s$ such that
$d_n\ge\delta$ for all $n\in s$,
\[
|s|\delta
\le
\frac1{\sqrt{c_*}}
\sqrt{\sum_n\varepsilon_n}\sqrt{\sum_n\tau_n},
\]
and consequently $\{n:d_n\ge\delta\}$ is finite. The exact manuscript-form
Lean declarations are listed in the source file as the counting theorem and
its corresponding finiteness theorem.

\paragraph{A.7 The finite-state operational bridge (Section~\ref{sec:metrics}).}
For a finite state space and a stochastic measurement channel $\Phi$,
\[
d_{\rm TV}(\Phi p,\Phi q)\le d_{\rm TV}(p,q).
\]
Hence any operational distance defined as the supremum of readout total
variation over a fixed admissible channel class satisfies
$d_{\rm op}\le d_{\rm TV}$. The Lean file
\nolinkurl{OperationalBridge.lean} verifies the $L^1$ contraction,
the total-variation data-processing inequality, sharpness for the identity
channel, transfer of a total-variation speed-limit hypothesis to
$d_{\rm op}$, and the fixed-resolution counting bound
\[
K_\delta^{\rm op}
\le
\frac1\delta\sqrt{\frac{\Sigma_*\mathcal N_*}{2}}.
\]
The stochastic-thermodynamic premise
$2d_{\rm TV}^2\le\Sigma\mathcal N$ is intentionally a hypothesis of the
Lean development rather than a theorem of it.

\subsection{What is not verified}

The following remain physical or modelling claims and are not certified by Lean:

\begin{enumerate}
\item the applicability of $\varepsilon_n\tau_n\ge c_nd_{\rm phys,n}^2$ to a
particular dynamics;
\item a universal identification or comparison of arbitrary physical metrics
with an operational organizational distance; the finite-state
total-variation/channel comparison of A.7 is machine-checked separately;
\item the claim that $\varepsilon_n$ is the correct cumulatively charged cost;
\item the stochastic-process theorem connecting the reciprocal-series criterion
to pure-birth explosion or non-explosion;
\item the physical additivity assumption $m_{\min}N_n\le M_n$.
\end{enumerate}

Thus this appendix does not show that the physical paper as a whole is formally
proved. It shows that the algebraic and analytic implications formalized here
are free of derivational gaps; the remaining burden lies in the explicitly
stated physical and stochastic-process premises.

\subsection{Reproduction}

The checked sources are
\nolinkurl{OrganizationalDepth.lean} and
\nolinkurl{OperationalBridge.lean}. They use Lean~4.33.0 and
\texttt{Mathlib} commit \texttt{db584cd} (10 August 2026). From the
verification directory, the check is reproduced with
\begin{verbatim}
lake update
lake exe cache get
lake build
\end{verbatim}
This exact three-command route is exercised by the repository CI. Both sources
compile with no errors and contain no \texttt{sorry}. The axiom audit for all
manuscript-facing declarations reports only
\texttt{propext}, \texttt{Classical.choice}, and \texttt{Quot.sound}.

The manuscript-facing declarations in \nolinkurl{OrganizationalDepth.lean}
cover stochastic-route closure, the weighted and uniform finite-action
bounds, and the fixed-resolution counting/finiteness results. The
declarations in \nolinkurl{OperationalBridge.lean} cover total-variation
data processing, identity-channel sharpness, transfer of a TV speed-limit
hypothesis to operational distance, and the operational fixed-resolution
counting bound.

\endgroup


\end{document}